\section{The Normal One：普通人叙事背后的研究选择}

访谈开场从纽约、播客、童年和上海交大讲起。谢赛宁反复说自己不是 the chosen one，而是 the normal one。这句话容易被误读成谦虚的姿态；放到整期访谈里，它更像一种研究者自我定位：不要把成功解释成命运赐予的天才脚本，而要解释成长期选择、人与组织、好奇心、失败和偶然共同作用的结果。

他讲童年时，两个元素反复出现：一是母亲带着他四处旅行，二是父亲书房里的大量书。旅行让他接触真实世界的多样性，阅读让他获得抽象世界的入口。九岁有电脑、接触游戏和互联网后，他第一次感到“内容”和表达的爆炸。这些细节在人物稿里是成长故事，在技术笔记里则对应后面世界模型的两个面向：世界需要亲自进入，抽象需要持续建模。

\begin{knowledgebox}{从人物经历读技术路线}
这期访谈里，个人经历不是背景花絮。旅行、书、游戏、互联网、电影、纽约街头、研究实习和创业组织，都在回答同一个问题：智能系统到底应该从哪里学习世界？如果只从文本学习，它得到的是人类压缩后的叙述；如果要理解物理世界，就需要重新面对连续信号、身体行动、环境反馈和真实生活。
\end{knowledgebox}

\subsection{ACM 班、通识与不过度竞争}

前面讲童年和互联网入口，是为了说明兴趣从来不是单点出现的；本节转到 ACM 班，是因为这里第一次把兴趣放进了制度化训练环境。问题随之变成：什么样的教育环境更容易保护长期研究品味？

谢赛宁讲上海交大 ACM 班时，特别强调宽松和通识。比如“学子讲坛”要求学生讲任何与课程无关的东西，可以是哲学、历史、社会或科学。这段内容的教学意义是：早期训练如果只围绕排名和题目优化，很容易得到强执行者；如果允许学生建立广阔兴趣，就更可能形成 research taste。

他也明确说不喜欢过度竞争。这里不是反对竞争，而是反对把所有人压成单一指标。研究的长期价值常常来自非线性路径：先有兴趣，后有问题；先有误打误撞的经历，后有能串起来的主线。过度竞争会缩短时间尺度，让人只追逐当前可计分的任务。

\begin{warningbox}{不要把“普通人”理解成低目标}
The normal one 不是降低目标，而是拒绝神话式自我叙事。它把注意力从“我是不是天选之子”转向“我能不能持续选择值得做的问题、找到值得共事的人、承担长期的不确定性”。
\end{warningbox}

\subsection{本章小结}

开场的成长故事为整期节目定调：谢赛宁关心的不是单点胜负，而是长期轨迹。这个轨迹后来表现为：不按最确定的名校/大厂路径走，反复选择更贴近自己问题意识的人和组织。

